\documentclass[11pt]{article}
\usepackage[margin=1in]{geometry}
\usepackage{amsmath,amssymb}
\usepackage{enumitem}
\usepackage{booktabs}
\usepackage{hyperref}

\newcommand{\D}[2]{D_{#1}^{#2}}
\newcommand{\Ec}[2]{\mathbb{E}^{\mathbb{Q}_{#1}}_{t}\!\left[#2\right]}

\title{FRE-GY 9743, Assignment 6, Part I\\[2pt]\large An EUR/USD Forward
Collateralised in GBP}
\author{Sumeet Mishra}
\date{\today}

\begin{document}
\maketitle

\section*{Notation}

The assignment's $\mathrm{df}^{[\cdot,\cdot]}$ is a pair, so to keep the derivation
unambiguous define
\[
  \D{c}{y}(T) \;=\;
  \begin{minipage}[t]{0.68\textwidth}
  today's value, in currency $y$ at spot, of one unit of currency $y$ paid at $T$,
  when the trade is fully collateralised in currency $c$.
  \end{minipage}
\]
The self-collateralised case is the ordinary OIS curve:
$\D{\mathrm{GBP}}{\mathrm{GBP}} = \mathrm{df}^{\mathrm{SONIA}}$,
$\D{\mathrm{EUR}}{\mathrm{EUR}} = \mathrm{df}^{\mathrm{EuroSTR}}$,
$\D{\mathrm{USD}}{\mathrm{USD}} = \mathrm{df}^{\mathrm{SOFR}}$.

Section~4 maps this onto the assignment's notation, under both readings of the pair.
$S(0)$ is EUR/USD spot in the market convention of USD per EUR, and $X_{a\to b}$ is the
rate converting one unit of $a$ into $b$.

\section{Q1. The valuation formula}

\subsection{Framework}

Under a CSA that is fully collateralised in currency $c$ with the collateral remunerated
at that currency's overnight rate $c_t$, the collateral account is the numeraire, and the
value in currency $c$ of any payoff is
\begin{equation}
  V^c(t) = \Ec{c}{e^{-\int_t^T c_s\,\mathrm{d}s}\,V^c(T)}.
  \label{eq:csa}
\end{equation}
This is the point of Assignment 4: posting collateral is not settlement, the trade keeps
its value, and what the collateral earns decides the discount rate.

Apply \eqref{eq:csa} to the single cashflow ``one unit of currency $y$ at $T$'', whose
value at $T$ in currency $c$ is $X_{y\to c}(T)$:
\begin{equation}
  \Ec{c}{e^{-\int_t^T c_s\,\mathrm{d}s}\,X_{y\to c}(T)}
  \;=\; X_{y\to c}(t)\,\D{c}{y}(t,T),
  \label{eq:defD}
\end{equation}
which is the definition of $\D{c}{y}$ written as a pricing statement. Note the asymmetry
that drives everything below: the \emph{cashflow} currency is $y$ but the \emph{discount}
is governed by $c$, so $\D{c}{y}$ is a genuinely two-index object. Only when $c=y$ does it
collapse to a plain OIS curve.

\subsection{The forward}

Long $N$ EUR of an EUR/USD forward at strike $K$: at $T$ receive $N$ EUR and pay $NK$
USD. Collateral is GBP, so $c = \mathrm{GBP}$ throughout. Taking GBP as the valuation
currency and applying \eqref{eq:defD} to each leg,
\[
  V^{\mathrm{GBP}}(0)
  = N\,X_{\mathrm{EUR}\to\mathrm{GBP}}(0)\,\D{\mathrm{GBP}}{\mathrm{EUR}}(T)
  \;-\; N K\,X_{\mathrm{USD}\to\mathrm{GBP}}(0)\,\D{\mathrm{GBP}}{\mathrm{USD}}(T).
\]
Dividing by $X_{\mathrm{USD}\to\mathrm{GBP}}(0)$ to express the result in USD, and using
$X_{\mathrm{EUR}\to\mathrm{GBP}}/X_{\mathrm{USD}\to\mathrm{GBP}} = X_{\mathrm{EUR}\to\mathrm{USD}} = S$,
\begin{equation}
  \boxed{\;
  V^{\mathrm{USD}}(0)
  = N\Bigl[\,S(0)\,\D{\mathrm{GBP}}{\mathrm{EUR}}(T)
  \;-\; K\,\D{\mathrm{GBP}}{\mathrm{USD}}(T)\Bigr].\;}
  \label{eq:value}
\end{equation}

Two observations worth stating explicitly, because they are what the question is probing:

\begin{itemize}[leftmargin=*]
\item \textbf{The collateral currency enters only through the discount factors, and it
  enters both legs.} It changes neither cashflow. A GBP CSA does not make this a GBP
  trade.
\item \textbf{No GBP curve appears.} $\D{\mathrm{GBP}}{\mathrm{GBP}}$ is absent from
  \eqref{eq:value} because the trade has no GBP cashflow. Posting GBP collateral means
  GBP governs \emph{how} the EUR and USD legs are discounted; it does not introduce a GBP
  payment to discount. This is why $\mathrm{df}^{\mathrm{SONIA}}$ turns out not to be
  needed in Section~3. That looks surprising until the asymmetry in
  \eqref{eq:defD} is taken seriously.
\end{itemize}

\section{Q2. The fair forward strike}

The fair strike is the $K$ that makes \eqref{eq:value} vanish:
\begin{equation}
  \boxed{\;
  K^\ast = S(0)\,
  \frac{\D{\mathrm{GBP}}{\mathrm{EUR}}(T)}{\D{\mathrm{GBP}}{\mathrm{USD}}(T)}.\;}
  \label{eq:strike}
\end{equation}

Substituting \eqref{eq:strike} back into \eqref{eq:value} gives the form actually used to
mark an existing trade:
\begin{equation}
  V^{\mathrm{USD}}(0) = N\,\D{\mathrm{GBP}}{\mathrm{USD}}(T)\,\bigl(K^\ast - K\bigr).
  \label{eq:mtm}
\end{equation}
The split in \eqref{eq:mtm} matters for Section~3: the \emph{strike} depends on the
\emph{ratio} of the two GBP-collateralised curves, while the \emph{value} additionally
needs one of them on its own.

\section{Q3. Are the five discount factors sufficient?}

\subsection{What is needed, and what is missing}

Equations \eqref{eq:value} and \eqref{eq:strike} require exactly two objects:
$\D{\mathrm{GBP}}{\mathrm{EUR}}$ and $\D{\mathrm{GBP}}{\mathrm{USD}}$. Of these,
$\D{\mathrm{GBP}}{\mathrm{USD}}$ is supplied directly. \textbf{Strictly,
$\D{\mathrm{GBP}}{\mathrm{EUR}}$ is not in the list}. The EUR curve under GBP
collateral is precisely the one that is absent.

So the literal answer is: not sufficient as a matter of pure bookkeeping. The useful
answer is that the missing curve is \emph{constructible} from the rest, provided one
additional assumption is stated rather than smuggled in.

\subsection{The assumption that closes the gap}

The relevant assumption is that a fully collateralised FX forward has the same fair rate
regardless of which currency the CSA specifies. Equivalently: for any pair $(a,b)$ and any
collateral currency $c$,
\begin{equation}
  F_{a\to b}(T) = X_{a\to b}(0)\,\frac{\D{c}{a}(T)}{\D{c}{b}(T)}
  \quad\text{is independent of } c.
  \label{eq:consistency}
\end{equation}
In basis language this is triangular consistency of the cross-currency bases: writing
$\D{c}{y} = \D{y}{y}\exp\bigl(-\!\int b_{c,y}\bigr)$, \eqref{eq:consistency} is the
statement $b_{\mathrm{GBP},\mathrm{EUR}} = b_{\mathrm{GBP},\mathrm{USD}} -
b_{\mathrm{EUR},\mathrm{USD}}$. It is the assumption every curve desk makes when building
a non-USD-collateral basis, because USD is the hub currency where cross-currency swap
liquidity actually sits.

Applying \eqref{eq:consistency} to the EUR/USD pair with $c = \mathrm{GBP}$ and
$c = \mathrm{EUR}$:
\[
  \frac{\D{\mathrm{GBP}}{\mathrm{EUR}}}{\D{\mathrm{GBP}}{\mathrm{USD}}}
  = \frac{\D{\mathrm{EUR}}{\mathrm{EUR}}}{\D{\mathrm{EUR}}{\mathrm{USD}}}
  \qquad\Longrightarrow\qquad
  \D{\mathrm{GBP}}{\mathrm{EUR}}
  = \D{\mathrm{GBP}}{\mathrm{USD}}\;
    \frac{\D{\mathrm{EUR}}{\mathrm{EUR}}}{\D{\mathrm{EUR}}{\mathrm{USD}}},
\]
and every term on the right is supplied. \textbf{So the answer is yes: sufficient,
conditional on \eqref{eq:consistency}.}

\subsection{What the formulas collapse to}

Substituting into \eqref{eq:strike}, the GBP factor cancels:
\begin{equation}
  K^\ast
  = S(0)\,\frac{\D{\mathrm{EUR}}{\mathrm{EUR}}(T)}{\D{\mathrm{EUR}}{\mathrm{USD}}(T)}
  = S(0)\,
  \frac{\mathrm{df}^{\mathrm{EuroSTR}}(T)}{\mathrm{df}^{[\mathrm{EUR},\mathrm{USD}]}(T)} .
  \label{eq:strike_final}
\end{equation}

This is the whole content of the question in one line. \textbf{The fair strike does not
depend on the collateral currency at all.} That is not a coincidence or a lucky
cancellation. It is \eqref{eq:consistency} restated, since that assumption says exactly
that the forward rate is collateral-invariant. A GBP CSA buys you nothing and costs you
nothing on a trade struck at the market.

The value of an \emph{off-market} forward is a different matter. From \eqref{eq:mtm},
\begin{equation}
  V^{\mathrm{USD}}(0)
  = N\,\mathrm{df}^{[\mathrm{GBP},\mathrm{USD}]}(T)\,(K^\ast - K),
  \label{eq:mtm_final}
\end{equation}
and here the GBP collateralisation genuinely bites: the payoff $(K^\ast - K)$ is discounted
on the GBP-collateralised USD curve, not on SOFR. Two desks holding the same trade under
different CSAs agree on the fair strike and disagree on the mark.

\subsection{Which of the five are actually used}

\begin{center}
\begin{tabular}{llc}
\toprule
Curve & Role & Needed? \\
\midrule
$\mathrm{df}^{\mathrm{EuroSTR}}$          & numerator of the strike \eqref{eq:strike_final}      & yes \\
$\mathrm{df}^{[\mathrm{EUR},\mathrm{USD}]}$ & denominator of the strike \eqref{eq:strike_final}   & yes \\
$\mathrm{df}^{[\mathrm{GBP},\mathrm{USD}]}$ & discounting of the mark \eqref{eq:mtm_final}        & yes \\
$\mathrm{df}^{\mathrm{SOFR}}$             & USD under \emph{USD} collateral; wrong CSA here      & no \\
$\mathrm{df}^{\mathrm{SONIA}}$            & GBP cashflows; the trade has none                    & no \\
\bottomrule
\end{tabular}
\end{center}

Two of the five are redundant for this trade, and the two that are redundant are the ones
a careless answer reaches for first: SONIA because the collateral is GBP, SOFR because the
trade settles in USD. Neither reason survives \eqref{eq:defD}. SOFR would be needed if the
CSA were in USD; SONIA only if the trade paid GBP.

\subsection{Mapping to the assignment's notation}

The pair $\mathrm{df}^{[\cdot,\cdot]}$ admits two readings, and the structural answer is
the same either way; only the arithmetic route changes.

\begin{description}[leftmargin=*,style=nextline]
\item[Reading A, $\mathrm{df}^{[c,y]}$ is cashflow $y$ under collateral $c$
  (used above).]
  $\mathrm{df}^{[\mathrm{GBP},\mathrm{USD}]} = \D{\mathrm{GBP}}{\mathrm{USD}}$ is given
  directly, and $\D{\mathrm{GBP}}{\mathrm{EUR}}$ is built as above. SONIA and SOFR are
  unused.
\item[Reading B, $\mathrm{df}^{[y,c]}$ is cashflow $y$ under collateral $c$.]
  Now $\mathrm{df}^{[\mathrm{EUR},\mathrm{USD}]} = \D{\mathrm{USD}}{\mathrm{EUR}}$ and
  $\mathrm{df}^{[\mathrm{GBP},\mathrm{USD}]} = \D{\mathrm{USD}}{\mathrm{GBP}}$. Applying
  \eqref{eq:consistency} to EUR/USD gives the strike straight off USD-collateral curves,
  $K^\ast = S(0)\,\mathrm{df}^{[\mathrm{EUR},\mathrm{USD}]}/\mathrm{df}^{\mathrm{SOFR}}$,
  and applying it to the USD/GBP pair gives the discounting curve
  $\D{\mathrm{GBP}}{\mathrm{USD}} = \mathrm{df}^{\mathrm{SONIA}}\,
  \mathrm{df}^{\mathrm{SOFR}}/\mathrm{df}^{[\mathrm{GBP},\mathrm{USD}]}$. Here four of
  the five are used and EuroSTR is the redundant one.
\end{description}

Under either reading: sufficient given \eqref{eq:consistency}, the strike is
collateral-independent, and the GBP CSA shows up only in the discounting of the mark.

\subsection{When the assumption fails}

Equation \eqref{eq:consistency} is a modelling choice, not a theorem, and it is worth
being precise about where it breaks:

\begin{enumerate}[label=\arabic*.,leftmargin=*]
\item \textbf{A directly quoted EUR/GBP basis.} The triangulated
  $\D{\mathrm{GBP}}{\mathrm{EUR}}$ is only as good as the claim that EUR/GBP basis equals
  GBP/USD minus EUR/USD. EUR/GBP cross-currency swaps do trade, and when the quoted basis
  differs from the triangulated one, the constructed curve is simply wrong. If that market
  is liquid at the relevant tenor, build $\D{\mathrm{GBP}}{\mathrm{EUR}}$ from it directly
  and treat the identity as a check rather than as the source.
\item \textbf{Collateral optionality.} A multi-currency CSA that lets the poster choose
  which currency to deliver makes the discount curve an option on the cheapest-to-deliver
  basis, not a curve at all. Then no set of single-currency discount factors is
  sufficient, however large. The question specifies GBP collateral, which avoids this,
  but it is the first thing to confirm in the actual CSA, because the phrase ``collateralised
  in GBP'' is often shorthand for ``GBP is one of the eligible currencies''.
\item \textbf{Collateral convexity.} \eqref{eq:consistency} holds exactly under
  deterministic rates and FX. With both stochastic, correlation between the collateral
  rate and the FX rate puts a convexity correction between the collateral-adjusted
  curves. It is second-order and routinely ignored for short tenors, which is the honest
  reason the identity is used rather than any claim that it is exact.
\end{enumerate}

\end{document}
